\chapter{通常の対称性からグラスマン代数へ}

\lectureoverview{超対称性は対称性であるから、まず量子力学においてこの語が何を意味するかを思い出そう。物理的な変換はヒルベルト空間上のユニタリ演算子によって表現され、連続的な変換は無限小に生成され、その交換子が変換の合成の仕方を符号化する。真の対称性はエネルギーを保存し、これはハミルトニアンと交換することと等価である。超対称性は、よく知られたすべての例が共有する一つの仮定を変える。すなわち、その生成子はボソンをフェルミオンに、フェルミオンをボソンに変える。そのような生成子を構成することは、自然に生成消滅演算子、反交換子、そして奇な量の二乗がゼロになる有限のグラスマン代数へと導く。}

\section{対称性は異なる配置を関係づける}

対称性は単なる装飾的な規則性ではない。それは一つの物理的配置を別の配置に関係づける操作である。例えば左右反転は両手を関係づける。もし人と周囲の法則が厳密な左右対称性をもつなら、両手の対応する性質は一致するであろう。より一般に、厳密な対称性によって結ばれた対象や状態は、その対称性が保存する物理量を共有しなければならない。素粒子物理学ではそれはしばしば等しい質量を意味し、原子物理学では同じ回転多重項に属する状態が等しいエネルギーをもつことを意味する。

我々の目的は超対称性であるが、まず回転や並進のような連続対称性を思い出すのが有用である。超対称性が通常の意味で「連続的」であるかという問いは、すでに厄介である。その無限小パラメータは通常の可換な実数ではない。それは連続対称性の次数付き拡張とみなすのが最善であり、以下で導入するグラスマン演算が、その主張を我々の当面の目的に十分なほど正確にする。\editorialnote{現代の用語では、超対称性を連続的な次数付き対称性として記述する。この限定は、通常の連続対離散という問いが超対称性にきれいに当てはまらないという講義の警告を説明するものである。}

\section{ヒルベルト空間上のユニタリ操作}

量子系は状態ベクトル\(\ket{\psi}\)によって表される。物理的な対称性操作は、その状態に作用する演算子\(U\)によって表される：
\begin{equation}
\ket{\psi}\longmapsto U\ket{\psi}.
\end{equation}
空間回転の場合、軸と角度が\(U\)を決定する。状態がいくつかの位置に粒子を含むなら、変換されたケットはそれらすべてを回転後の位置に表す。並進も同様に働く。すなわち、配置全体が変位した状態に置き換える。

\begin{classroomqa}
\audiencequestion{状態ベクトルは、しばしば無限次元のヒルベルト空間に属する。我々はヒルベルト空間を回転しているのか、それとも通常の三次元空間を回転しているのか。\lecturetimestamp{00:04:20}}
\lecturerresponse{どちらの記述も、異なるレベルで同じ操作を指している。ユニタリ変換はヒルベルト空間のベクトルに作用し、通常の空間における配置の物理的回転を表す。それは抽象的なヒルベルト空間の文字通りの三次元回転ではなく、空間回転のそこでの表現である。}
\end{classroomqa}

本質的な条件はユニタリ性である：
\begin{equation}
U^\dagger U=UU^\dagger=\mathbf 1,
\qquad
U^{-1}=U^\dagger.
\end{equation}
これは内積を、したがって確率を保存する：
\begin{equation}
\braket{U\psi}{U\psi}
=\bra{\psi}U^\dagger U\ket{\psi}
=\braket{\psi}{\psi}.
\end{equation}

\begin{figure}[H]
\centering
\includegraphics[width=0.56\textwidth]{lecture_04_figure_02.png}
\lectureframe{00:05:56.420}
\caption{初期の板書では、変換された状態からユニタリ性の条件へと移行する。ケットとダガー付きの積は見えるが、講師がそれらの右端を隠している。完全な式は以下の導出に現れる。}
\end{figure}

厳密に言えば、ウィグナーの定理は時間反転のような反ユニタリ対称性も許す。ここで扱う、恒等変換に連続的につながる変換はユニタリに表現される。\editorialnote{反ユニタリの可能性は標準的な限定であり、回転、並進、または超対称性生成子に関する講義の議論を変えるものではない。}

\section{無限小変換と生成子}

何もしないことは\(U=\mathbf 1\)によって表され、\(U=0\)ではない。十分小さい変換は、その小さなパラメータに線形な項だけ恒等変換と異なる。小さな角度\(\epsilon\)の回転に対して、
\begin{equation}
U(\epsilon)=\mathbf 1+i\epsilon L+O(\epsilon^2).
\end{equation}
一次のユニタリ性は
\begin{align}
U^\dagger U
&=\left(\mathbf 1-i\epsilon L^\dagger\right)
  \left(\mathbf 1+i\epsilon L\right)+O(\epsilon^2)\\
&=\mathbf 1+i\epsilon(L-L^\dagger)+O(\epsilon^2),
\end{align}
を要求し、したがって
\begin{equation}
L^\dagger=L.
\end{equation}
演算子\(L\)は無限小生成子である。

回転には三つの生成子が必要であり、各独立な軸に一つずつ対応する：
\begin{equation}
L_x,\qquad L_y,\qquad L_z.
\end{equation}
任意の軸\(\widehat{\bm n}\)のまわりの回転は
\begin{equation}
L_{\widehat{\bm n}}=\widehat{\bm n}\cdot\bm L.
\end{equation}
によって生成される。

\begin{classroomqa}
\audiencequestion{\(x\)軸、\(y\)軸、\(z\)軸の生成子を足し合わせて、別の軸のまわりの回転をつくることはできますか。\lecturetimestamp{00:07:50}}
\lecturerresponse{はい。それらの線形結合は中間的な軸のまわりの回転を生成します。3つの独立な成分があれば、あらゆる無限小の空間回転を構成するのに十分です。}
\end{classroomqa}

回転の場合、これらの生成子は角運動量の成分である。この記法はより一般に有用である。なぜなら、あらゆる連続対称性はその生成子を通じて局所的に理解されるからである。

\section{乗積表としての交換子}

回転の生成子は
\begin{align}
[L_x,L_y]&=iL_z,&
[L_y,L_z]&=iL_x,&
[L_z,L_x]&=iL_y,
\end{align}
を満たす。あるいは簡潔に
\begin{equation}
[L_i,L_j]=i\epsilon_{ijk}L_k.
\end{equation}
と書ける。Levi-Civita記号は、2つの添字が一致すれば消え、\(1,2,3\)の偶置換に対しては\(+1\)、奇置換に対しては\(-1\)である。重要な構造上の事実は閉性である。すなわち、2つの生成子の交換子は再び生成子になる。

なぜ交換子が回転の幾何学を記述するのだろうか。\(x\)のまわりの小さな回転を行い、次に\(y\)のまわりの小さな回転を行い、\(x\)回転を元に戻し、最後に\(y\)回転を元に戻す。異なる軸のまわりの回転は可換でないため、これらの操作は打ち消し合わない。

\begin{classroomqa}
\audiencequestion{この非可換性は、飛行機のロール、ピッチ、ヨーを通じて描くことができますか。\lecturetimestamp{00:12:36}}
\lecturerresponse{はい。機体軸のまわりの逐次回転は、順序を変えると最終的な向きが変わることを即座に幾何学的に示します。以下の代数は、そのよく知られた事実の無限小版です。}
\end{classroomqa}

ここで
\begin{equation}
U_x(\epsilon)=\mathbf 1+i\epsilon L_x,
\qquad
U_y(\delta)=\mathbf 1+i\delta L_y.
\end{equation}
とする。演算子はケットに右から左へ作用するので、述べられた順序は
\begin{equation}
\mathcal C
=U_y^{-1}(\delta)U_x^{-1}(\epsilon)
 U_y(\delta)U_x(\epsilon).
\end{equation}
で表される。

\begin{classroomqa}
\audiencequestion{なぜ元に戻す操作はダガーではなく逆元として書けるのですか。\lecturetimestamp{00:13:25}}
\lecturerresponse{変換がユニタリだからです。すなわち\(U^\dagger=U^{-1}\)です。無限小変換では、逆元を取ることは1次の項の符号を反転させるだけです。\(U_x^{-1}=\mathbf 1-i\epsilon L_x+O(\epsilon^2)\)}
\end{classroomqa}

黒板で4つの因子を掛け合わせることは活発な練習になった。\(\epsilon\)または\(\delta\)について線形な項はすべて打ち消し合い、教室中が混合項を追い、一時的な因子2の混乱は演算子の順序を失うことがいかに容易かを露呈した。きれいな2次の結果は
\begin{align}
\mathcal C
&=\mathbf 1+\epsilon\delta[L_x,L_y]
  +O(\epsilon^2\delta,\epsilon\delta^2)\\
&=\mathbf 1+i\epsilon\delta L_z
  +O(\epsilon^2\delta,\epsilon\delta^2).
\end{align}
である。したがって、この閉じた順序は\(\epsilon\delta\)に比例する角度だけ\(z\)のまわりの小さな回転を残す。\editorialnote{表示された掛け算は、教室での探索的な計算の後に順序と係数を確定する。順序についての規約を逆にすると符号は逆になるが、残差変換が交換子によって支配されるという結論は変わらない。}

\begin{classroomqa}
\audiencequestion{逆操作の順序を変えると積は恒等演算子になりますか。\lecturetimestamp{00:25:23}}
\lecturerresponse{各変換をそれ自身の逆元のすぐ隣に置けば、対は打ち消し合います。非自明な交換子の順序は意図的にそれらの対を引き離しており、それが潰れないことが順序依存性を測っています。}
\end{classroomqa}

より長い積は確かにより複雑になるが、そこに新しい原理は含まれていない。

\begin{classroomqa}
\audiencequestion{生成子が3つより多い場合、あるいは小さな変換が多数ある場合はどうなりますか。\lecturetimestamp{00:25:34}}
\lecturerresponse{交換子が増え、記帳も増えますが、閉じた代数を知っていれば、任意に多くの無限小変換がどのように合成されるかを決定できます。有限変換は多くの小さな変換から組み立てることができます。}
\end{classroomqa}

この意味で、交換子代数は恒等演算子の近くで書かれた連続群の乗積表である。

\section{並進が空間の例を完成させる}

空間並進は運動量によって生成される。講義で用いた符号の規約では、
\begin{equation}
\begin{aligned}
T_x(a)&=\mathbf 1+iaP_x+O(a^2),\\
T_y(b)&=\mathbf 1+ibP_y+O(b^2),
\end{aligned}
\end{equation}
である。ここで\(a\)と\(b\)は角度ではなく距離である。

2つの並進は可換である。点を\(x\)方向に変位させ、次に\(y\)方向に変位させ、\(x\)並進を元に戻し、\(y\)並進を元に戻すと、点は元の位置に戻る。
\begin{equation}
[P_i,P_j]=0.
\end{equation}

\begin{classroomqa}
\audiencequestion{物理的な物体を並進または回転させるには、運動を始めたり止めたりする力が必要ですか。\lecturetimestamp{00:28:45}}
\lecturerresponse{それはここで定義している操作ではありません。対称変換は、元の状態と、配置全体が変位または回転された数学的な状態とを比較します。それは力学学的な軌道、加速度、あるいは装置を動かす間になされる仕事を記述するものではありません。}
\end{classroomqa}

回転と並進は可換ではない。点を並進させ、原点を中心に回転させ、元の軸に沿って逆に並進させ、回転を元に戻すと、点は最初の位置には戻らない。混合代数は
\begin{equation}
[L_i,P_j]=i\epsilon_{ijk}P_k.
\end{equation}
である。この関係は、運動量自体が回転の下でベクトルとして変換することを述べている。ここでは厳密な添字の練習よりも、一般的な教訓の方が重要である。すなわち、ユークリッド空間の変換は、その回転生成子と並進生成子の間の交換子に符号化されている。

\section{対称性はエネルギーを保存しなければならない}

ここで変換の合成から力学へ移る。\(U\) が真に系の対称性であるならば、それを適用してもエネルギーは変化しえない。\(\ket{E}\) をエネルギー固有状態とする：
\begin{equation}
H\ket{E}=E\ket{E}.
\end{equation}
右辺の記号 \(E\) は数であり、\(\ket{E}\) はその固有値でラベル付けされたベクトルである。

\begin{figure}[H]
\centering
\begin{minipage}[t]{0.43\textwidth}
\centering
\includegraphics[width=\linewidth]{lecture_04_figure_03.png}
\lectureframe{00:34:38.239}
\end{minipage}
\hfill
\begin{minipage}[t]{0.50\textwidth}
\centering
\includegraphics[width=\linewidth]{lecture_04_figure_04.png}
\lectureframe{00:36:22.376}
\end{minipage}
\caption{変換後のケット \(U\ket{E}\) と、定義する固有値方程式 \(H\ket{E}=E\ket{E}\) を、連続する板書のフレームに分離して示したもの。}
\end{figure}

\(U\) が対称性ならば、\(U\ket{E}\) は同じ固有値をもつ：
\begin{equation}
HU\ket{E}=EU\ket{E}.
\end{equation}
\(E\) はスカラーなので \(U\) と可換であり、元の固有値方程式を用いることができる：
\begin{align}
HU\ket{E}
&=EU\ket{E}\\
&=UE\ket{E}\\
&=UH\ket{E}.
\end{align}
したがって
\begin{equation}
(HU-UH)\ket{E}=0.
\end{equation}

\begin{classroomqa}
\audiencequestion{この導出は最初の関係から繰り返すことができるか？\lecturetimestamp{00:35:55}}
\lecturerresponse{\(H\ket{E}=E\ket{E}\) から始める。対称性により \(U\ket{E}\) は同じエネルギーをもつので、\(HU\ket{E}=EU\ket{E}\) である。スカラー \(E\) を \(U\) を通して移動し、\(E\ket{E}\) を \(H\ket{E}\) で置き換え、引き算する。結果は \((HU-UH)\ket{E}=0\) である。}
\end{classroomqa}

\begin{figure}[H]
\centering
\includegraphics[width=0.70\textwidth]{lecture_04_figure_05.png}
\lectureframe{00:38:51.066}
\caption{完全な板書による導出。曲線のガイドマークが等エネルギー文から置換 \(E\ket{E}=H\ket{E}\) へとつながり、最後の行は \((HU-UH)\ket{E}=0\) である。}
\end{figure}

この議論はすべてのエネルギー固有状態に対して成り立つ。それらの状態が完全な基底をなすと仮定すると、それらすべてを消す演算子はゼロ演算子である：
\begin{equation}
[H,U]=0.
\end{equation}
これが、時間に依存しないユニタリ変換がハミルトニアンの対称性であるための判定基準である。

無限小変換
\begin{equation}
U(\epsilon)=\mathbf 1+i\epsilon L_i+O(\epsilon^2),
\end{equation}
に対して、恒等演算子はすでに \(H\) と可換なので、各生成子は
\begin{equation}
[L_i,H]=0.
\end{equation}
を満たさなければならない。

\begin{figure}[H]
\centering
\includegraphics[width=0.47\textwidth]{lecture_04_figure_06.png}
\lectureframe{00:42:14.296}
\caption{板書に要約された対称性の判定基準：\([U,H]=0\)、続いてその生成子形。部分的に不鮮明な下付きラベルから \([L_i,H]=0\) として標準化。}
\end{figure}

\(L_i\) が陽に時間に依存しないならば、ハイゼンベルク方程式は
\begin{equation}
\frac{dL_i}{dt}
=i[H,L_i]=0
\qquad(\hbar=1).
\end{equation}
を与える。したがって対称性の生成子は保存電荷である。空間並進対称性は運動量保存を与え、回転対称性は角運動量保存を与える。\editorialnote{陽に時間微分するという条件は標準的であり、講義の保存則の記述に暗黙に含まれている。}

\section{物理系全体を変換せよ}

「回転は対称性である」という表現は、装置の一部だけを変換する場合には誤用されうる。自由なスピンはエネルギーを変えずに回転させることができる。しかし、それを固定した磁場の中に置き、スピンだけを回転させると、そのゼーマンエネルギーは変化する。これは回転対称性を反証するものではない。磁場を固定したまま、スピンをそれに対して回転させたことを意味する。スピンと磁場を一緒に回転させるか、それ自体が不変である外部磁場を用いれば、エネルギーは変化しない。

同じ問題は水素原子にも現れる。陽子に対して電子だけを並進させると、それらの間隔が変わり、したがってクーロンエネルギーが変わる。並進対称性は完全な電子-陽子系に作用し、両方を同じ量だけ動かす。

\begin{classroomqa}
\audiencequestion{外部場自体がその対称性をもたなければならないのか？\lecturetimestamp{00:43:09}}
\lecturerresponse{固定された外部場は、すでにその操作の下で不変であるか、さもなければ系の残りと一緒に変換されなければならない。スピンと回転されていない磁場の間の相対回転は、異なる物理的配置であり、系全体の対称性変換ではない。}
\end{classroomqa}

原子における回転縮退は明快な正の例を与える。固定された角運動量 \(l\) に対して、回転は磁気量子数
\begin{equation}
m=-l,-l+1,\ldots,l.
\end{equation}
をもつ状態を混合する。外部磁場または電場がなければ、多重項のすべての要素は同じエネルギーをもつ：
\begin{equation}
E_{l,m}=E_l.
\end{equation}
対称性は異なる状態を関係づけ、それらの縮退を強制する。

\section{超対称性が変えるもの}

これまでに論じてきたよく知られた対称性は、いずれもボソンとフェルミオンという性格を保つ。回転は電子を別の電子へ移すが、光子にはしない。カラーの回転はクォークの3つのカラーを混ぜるが、クォークをグルーオンに変えることはない。通常の内部対称性や時空対称性は成分や量子ラベルを混ぜることがあっても、その生成子はボソン的であり、フェルミオン・パリティを変えない。

超対称性は異なる：
\begin{equation}
Q\ket{B}\sim\ket{F},
\qquad
Q\ket{F}\sim\ket{B}.
\end{equation}
その多重項はボソンとフェルミオンの境界を越える。厳密な超対称理論では、同じ多重項に属するパートナーは同じ質量をもつ。これはまさに前の講義で要求された対合である。対応するボソン的ループとフェルミオン的ループが相殺して、スカラー質量への大きな紫外補正を消すことができる。通常の空間的な運動によって整数スピンの状態を半整数スピンの状態に変えることはできないため、この構成を視覚化するのは難しい。したがって、われわれは代数的に進める。

\section{超電荷の原型}

\(A,A^\dagger\) をボソンを消滅・生成する演算子、\(C,C^\dagger\) をフェルミオンを消滅・生成する演算子とする。それらの1粒子状態を入れ替える最も単純なエルミート演算子は
\begin{equation}
Q=C^\dagger A+A^\dagger C.
\end{equation}
である。第1項はボソンを消してフェルミオンを生成し、第2項はフェルミオンを消してボソンを生成する。\(Q\) が1ボソン状態に作用すると、消滅すべきフェルミオンが存在しないため \(A^\dagger C\) の項は消え、\(C^\dagger A\) が変換を行う。1フェルミオン状態では役割が逆になる。

\begin{classroomqa}
\audiencequestion{この構成は、ボソンとフェルミオンの重ね合わせである単一の粒子を生成するのですか？ \lecturetimestamp{00:52:37}}
\lecturerresponse{\(Q\) は状態ではなく演算子である。純粋なボソン状態に作用すればフェルミオン状態を生成し、純粋なフェルミオン状態に作用すればボソン状態を生成する。すでに一般的な重ね合わせに線形に作用すれば、任意の線形演算子がそうであるように、別の重ね合わせを生成する。}
\end{classroomqa}

次の部屋では真空を試す。両方の消滅演算子が真空を消すので、
\begin{equation}
A\ket{0}=0,\qquad C\ket{0}=0,
\end{equation}
したがって
\begin{equation}
Q\ket{0}=0.
\end{equation}
が成り立つ。結果はゼロベクトルであり、真空の2つ目のコピーではない。

この原型は意図的に模式的なものである。スピン0のボソンはスピン1/2のフェルミオンのどちらの成分にも写像できるので、生成子自身がスピノル指数を必要とする：
\begin{equation}
Q_\alpha\ket{B}\sim\ket{F_\alpha}.
\end{equation}
講義で用いた簡略化された2成分描像では、
\begin{equation}
Q_\alpha\ket{F_\beta}\sim\delta_{\alpha\beta}\ket{B}.
\end{equation}

\begin{classroomqa}
\audiencequestion{スピン0のボソンが1つのスピン成分をもち、スピン1/2のフェルミオンが2つをもつとき、状態空間の次元はどのように一致するのですか？ \lecturetimestamp{00:55:02}}
\lecturerresponse{指数をもたない単一の \(Q\) では不十分である。スピノル成分 \(Q_\alpha\) があり、それぞれがフェルミオンの各スピン成分を生成することができる。超電荷自身はスピンを半単位だけ変えるので、半単位のスピンをもつ。}
\end{classroomqa}

\(Q_\alpha\) はフェルミオン的、より正確にはグラスマン奇の演算子である。それらのすべての項は奇数個のフェルミオン演算子を含み、そのためにフェルミオン・パリティを反転させる。通常の可観測量は偶である。奇の超電荷とその奇の変換パラメータは区別すべきである。有限の超対称変換は、\(\epsilon^\alpha Q_\alpha+\bar\epsilon_{\dot\alpha}\bar Q^{\dot\alpha}\) のような偶の組み合わせから構成され、ここで \(\epsilon\) と \(Q\) はともに奇である。\editorialnote{講義では非公式に生成子をグラスマン数と呼んでいる。より正確には、超電荷は奇の演算子であり、反可換なパラメータはグラスマン数である。それらの積は偶であり、有限変換に現れることができる。}

\section{ボソン的演算子代数とフェルミオン的演算子代数}

\(i,j\) でラベル付けられたボソン的モードに対して、
\begin{align}
[A_i,A_j]&=0,&
[A_i^\dagger,A_j^\dagger]&=0,&
[A_i,A_j^\dagger]&=\delta_{ij}.
\end{align}
ラベルは運動量、位置、スピン、または任意の完全な1粒子基底を表してよい。生成演算子は可換なので、
\begin{equation}
A_i^\dagger A_j^\dagger\ket{0}
=A_j^\dagger A_i^\dagger\ket{0}.
\end{equation}
2ボソン状態は入れ替えに対して対称である。

フェルミオン演算子は反交換関係に従う：
\begin{align}
\{C_i,C_j\}&=0,&
\{C_i^\dagger,C_j^\dagger\}&=0,&
\{C_i,C_j^\dagger\}&=\delta_{ij},
\end{align}
ここで
\begin{equation}
\{X,Y\}=XY+YX.
\end{equation}
単一モードに対して、
\begin{equation}
\{C^\dagger,C^\dagger\}=2(C^\dagger)^2=0,
\end{equation}
したがって
\begin{equation}
(C^\dagger)^2=0.
\end{equation}
演算子はゼロではない；その二乗がゼロである。それはそのモードに一つのフェルミオンを生成できるが、二度目の適用はゼロを与える。これが演算子の言葉によるパウリ排他である。

異なるラベルに対して、
\begin{equation}
C_i^\dagger C_j^\dagger\ket{0}
=-C_j^\dagger C_i^\dagger\ket{0},
\end{equation}
これは2フェルミオン状態の反対称性である。代数的な置き換え
\begin{equation}
\begin{gathered}
\text{bosons: commutators}\\
\Updownarrow\\
\text{fermions: anticommutators}
\end{gathered}
\end{equation}
が統計的な区別を符号化している。

多数のボソンが1つのモードを占めてよい。したがって、大きな光子占有数は通常の数値振幅をもつ古典電磁場として振る舞うことができる。フェルミオンは1つのモードに蓄積できず、奇のフェルミオン場は通常の可換な古典値をもたない。このことがフェルミオン的記帳のための異なる算術を動機づける。

\section{グラスマン数とパリティ}

通常の数は可換である：
\begin{equation}
\alpha_i\alpha_j=\alpha_j\alpha_i.
\end{equation}
グラスマン生成元は反可換である：
\begin{equation}
\theta_i\theta_j=-\theta_j\theta_i.
\end{equation}
\(i=j\) とし、\(2\neq0\) が成り立つ数体系の上で置くと、
\begin{equation}
\theta_i^2=0.
\end{equation}
\(\theta_i\) に六や七のような通常の数値を割り当てることはできない。これらの記号はフェルミオン的符号を追跡するために設計された代数をなす。

通常の係数はグラスマン生成元と可換である：
\begin{equation}
\alpha\theta_i=\theta_i\alpha.
\end{equation}

\begin{classroomqa}
\audiencequestion{通常の数がグラスマン数に掛かると、どのような種類の対象が得られるか？ \lecturetimestamp{01:12:32}}
\lecturerresponse{それはグラスマン奇のままである。通常の係数による掛け算はそのスケールを変えるが、パリティは変えない。生成元の張る空間は通常の数の上のベクトル空間である。}
\end{classroomqa}

この代数は \(\mathbb Z_2\) パリティをもつ。通常の係数および偶数個の生成元の積は偶であり、生成元または奇数個の積は奇である：
\begin{equation}
\begin{aligned}
\text{even}\cdot\text{even}&=\text{even},\\
\text{even}\cdot\text{odd}&=\text{odd},\\
\text{odd}\cdot\text{odd}&=\text{even}.
\end{aligned}
\end{equation}

\begin{classroomqa}
\audiencequestion{これは実数と虚数の掛け算規則に類似しているのか、それとも \(i\) の導入のような一般化なのか？ \lecturetimestamp{01:14:25}}
\lecturerresponse{それは算術の一般化であるが、この類似を押し進めすぎてはならない。虚数は依然として可換で \(i^2=-1\) を満たすが、グラスマン奇の量は反可換で二乗するとゼロになる。偶奇の掛け算パターンを共有することは、それらの代数を同じものにはしない。}
\end{classroomqa}

\(\theta_1\theta_2\) のような積は偶であり、グラスマン代数のすべての斉次元と可換であるが、通常の数ではない。それは非零でありながら
\begin{equation}
(\theta_1\theta_2)^2=0.
\end{equation}

\begin{classroomqa}
\audiencequestion{\(\theta_1\theta_2\) が偶であるなら、なぜそれは単に通常の数ではないのか？ \lecturetimestamp{01:17:36}}
\lecturerresponse{偶パリティとは、それが次数付き代数において可換であることを意味するのであり、実数または複素スカラーであることを意味するのではない。その冪零性が、非零の通常の数とは直ちに区別する。}
\end{classroomqa}

\section{グラスマン変数の関数}

生成元が一つのとき、すべての多項式は一次で終わる：
\begin{equation}
f(\theta)=\alpha+\beta\theta.
\end{equation}
\(n\geq2\) であるすべての冪 \(\theta^n\) は消える。生成元が二つのとき、最も一般的な元は
\begin{equation}
f(\theta_1,\theta_2)
=\alpha+\beta_1\theta_1+\beta_2\theta_2
+\gamma\theta_1\theta_2.
\end{equation}
最後の単項式を逆順に書くと符号が変わる：
\begin{equation}
\theta_2\theta_1=-\theta_1\theta_2.
\end{equation}

\begin{classroomqa}
\audiencequestion{\(\theta_1\) に係数 \(\beta\) を掛けると、第三の独立なグラスマン変数が生じるか？ \lecturetimestamp{01:18:39}}
\lecturerresponse{いいえ。それは \(\theta_1\) の張る空間における別のベクトルを生じるのであり、新しい基底生成元を生じるのではない。二つの独立な生成元があるとき、すべての奇の線形項は \(\beta_1\theta_1+\beta_2\theta_2\) である。}
\end{classroomqa}

生成元が三つのとき、順序付けられた基底は
\begin{equation}
1,\quad
\theta_1,\theta_2,\theta_3,\quad
\theta_1\theta_2,\theta_1\theta_3,\theta_2\theta_3,\quad
\theta_1\theta_2\theta_3.
\end{equation}
一般に、\(n\) 個の生成元は \(2^n\) 個の順序付けられた基底単項式を生じる。有限個のグラスマン生成元の任意の解析関数は、生成元の繰り返しが消えるため、そのような多項式に切り詰められる。\editorialnote{\(2^n\) 次元の基底は、講義が一、二、三の生成元について構成するパターンを明示している。}

\section{グラスマン微分}

一変数のとき、
\begin{equation}
\frac{\partial}{\partial\theta}\alpha=0,
\qquad
\frac{\partial}{\partial\theta}\theta=1,
\end{equation}
と定義すると、
\begin{equation}
\frac{\partial}{\partial\theta}
(\alpha+\beta\theta)=\beta.
\end{equation}
二変数のとき、左微分を選び、標準的順序 \(\theta_1\theta_2\) を保つ：
\begin{align}
\frac{\partial^L f}{\partial\theta_1}
&=\beta_1+\gamma\theta_2,\\
\frac{\partial^L f}{\partial\theta_2}
&=\beta_2-\gamma\theta_1.
\end{align}
二行目のマイナス符号は、奇の微分が \(\theta_2\) を微分する前に \(\theta_1\) を通り抜けなければならないことから生じる。

これは次数付き積の規則によって表される。\(A\) がパリティ \(|A|\in\{0,1\}\) をもつなら、
\begin{equation}
\frac{\partial^L}{\partial\theta}(AB)
=\left(\frac{\partial^L A}{\partial\theta}\right)B
+(-1)^{|A|}A
\left(\frac{\partial^L B}{\partial\theta}\right).
\end{equation}
この規則は本質的には通常のものであるが、反可換性によって要求される符号を含んでいる。グラスマン積分も定義でき、関数空間が有限であるため、短い表に帰着する。講義ではその表を後回しにする。

\section{指数関数のパズル、解決}

単一のグラスマン生成元と通常の係数 \(a\) をもつとき、
\begin{equation}
e^{a\theta}
=1+a\theta,
\end{equation}
なぜならすべての高次の冪が消えるからである。同じ生成元を含む二つの指数関数は
\begin{align}
e^{a\theta}e^{b\theta}
&=(1+a\theta)(1+b\theta)\\
&=1+(a+b)\theta\\
&=e^{(a+b)\theta}.
\end{align}
を満たす。

次に教室では異なる生成元を試す。直接の掛け算は
\begin{equation}
e^{\theta_1}e^{\theta_2}
=(1+\theta_1)(1+\theta_2)
=1+\theta_1+\theta_2+\theta_1\theta_2.
\end{equation}
を与える。最初、黒板での展開は \(\theta_1+\theta_2\) の二乗が最後の項を与えるかもしれないことを示唆する。聴衆の一人が符号に気づく：
\begin{equation}
(\theta_1+\theta_2)^2
=\theta_1^2+\theta_1\theta_2
 +\theta_2\theta_1+\theta_2^2
=0.
\end{equation}

\begin{classroomqa}
\audiencequestion{二乗の項は \(\theta_1\theta_2+\theta_2\theta_1\) を含むために相殺するのか？ \lecturetimestamp{01:24:00}}
\lecturerresponse{はい。その修正は結論を変える：\(e^{\theta_1+\theta_2}=1+\theta_1+\theta_2\) である一方、順序付けられた積 \(e^{\theta_1}e^{\theta_2}\) も \(\theta_1\theta_2\) を含む。素朴な指数関数の加法法則は、異なる反可換な指数には適用されない。}
\end{classroomqa}

矛盾はない。通常の行列でさえ、\(A\) と \(B\) が可換であるときにのみ \(e^Ae^B=e^{A+B}\) を満たす。ここで
\begin{equation}
[\theta_1,\theta_2]
=\theta_1\theta_2-\theta_2\theta_1
=2\theta_1\theta_2.
\end{equation}
この交換子は偶であり、両方の生成子と可換であるから、Baker-Campbell-Hausdorff 級数は有限で終わる：
\begin{equation}
e^{\theta_1}e^{\theta_2}
=\exp\left(
\theta_1+\theta_2+\frac12[\theta_1,\theta_2]
\right)
=e^{\theta_1+\theta_2+\theta_1\theta_2}.
\end{equation}
右辺の指数は冪零であるから、ちょうど
\(1+\theta_1+\theta_2+\theta_1\theta_2\) に展開される。\editorialnote{この BCH 計算は、生講義が問いとして残した解決を与える。講義は素朴な法則の破綻を正しく見出したが、欠けている交換子項を特定する前に終わっている。}

\begin{classroomqa}
\audiencequestion{\(1+\theta_1\) はどのような数であり、通常部分と Grassmann 部分は加算してよいのか。\lecturetimestamp{01:26:16}}
\lecturerresponse{それは Grassmann 代数の非斉次元であり、偶部分 \(1\) と奇部分 \(\theta_1\) をもつ。そのような和は許されるが、通常のスカラーでも斉次な奇元でもない。}
\end{classroomqa}

終わり近くで、もう一つの見かけのパラドックスが検証される：
\begin{equation}
(\theta_1+\theta_2)^2=0.
\end{equation}
これは \(\theta_1\theta_2=0\) を意味しない。二つの交差項は別々に消えるのではなく相殺するからである。

\begin{classroomqa}
\audiencequestion{すべての奇 Grassmann 元が二乗してゼロになるなら、任意の二つの Grassmann 生成子の積はゼロであるということか。\lecturetimestamp{01:32:19}}
\lecturerresponse{いいえ。方程式 \((\theta_1+\theta_2)^2=0\) は \(\theta_1\theta_2+\theta_2\theta_1=0\) を述べている。それは反可換性を証明するのであって、どちらかの順序の積が消えることを証明するのではない。異なる生成子は非零の積をもちうる。}
\end{classroomqa}

\begin{classroomqa}
\audiencequestion{これらの指数関数を書くとき、通常の冪級数の定義を仮定しているのか。\lecturetimestamp{01:34:42}}
\lecturerresponse{はい、そしてその定義は各級数が有限で終わるため整合的である。破綻するのは指数関数そのものではなく、非可換な \(A\) と \(B\) に対する根拠のない恒等式 \(e^Ae^B=e^{A+B}\) である。交換子の補正が正しい合成則を回復する。}
\end{classroomqa}

この探索的な結びは適切である。通常の対称性は、ユニタリ変換から生成子と交換子へと我々を導いた。超対称性は、奇生成子から反可換変数へと我々を導いた。どちらの場合も、乗法則が中心的な対象である。回転についてはそれは \([L_i,L_j]\) に現れ、Grassmann 指数関数については、素朴な加法を妨げる交換子項に現れる。この代数は馴染みがないが、首尾一貫しており、有限的であり、超対称性が要求するフェルミオン的変換にまさに適合している。
